\documentclass{article}
\usepackage{amsmath,amssymb,amsthm}
\usepackage{algorithm}
\usepackage{algorithmic}

\newtheorem{theorem}{Theorem}
\newtheorem{lemma}[theorem]{Lemma}
\newtheorem{definition}[theorem]{Definition}

\begin{document}

\section*{Online Gradient Descent (OGD)}
\section*{Mathematical Formulation}
Let $\mathcal{W} \subseteq \mathbb{R}^d$ be a closed convex set. At each time step $t \in \{1, 2, \dots, T\}$, the algorithm chooses a point $w_t \in \mathcal{W}$. An adversarial then reveals a convex function $f_t : \mathcal{W} \to \mathbb{R}$. The algorithm incurs a loss $f_t(w_t)$ and computes the gradient $\nabla f_t(w_t)$. The update rule for Online Gradient Descent is:
\begin{align*}
w_{t+1} = \Pi_{\mathcal{W}} \left( w_t - \eta_t \nabla f_t(w_t) \right)
\end{align*}
where $\eta_t > 0$ is the learning rate (step size) at time $t$, and $\Pi_{\mathcal{W}}(w) = \arg\min_{v \in \mathcal{W}} \|v - w\|^2$ is the Euclidean projection onto the set $\mathcal{W}$.

\section*{Pseudocode}
\begin{algorithm}[H]
\caption{Online Gradient Descent (OGD)}
\begin{algorithmic}[1]
\REQUIRE Convex set $\mathcal{W}$, learning rate schedule $\{\eta_t\}_{t=1}^T$, iterations $T$
\STATE Initialize $w_1 \in \mathcal{W}$ arbitrarily
\FOR{$t = 1, 2, \dots, T$}
\STATE Play $w_t$
\STATE Observe loss function $f_t$ and incur loss $f_t(w_t)$
\STATE Compute gradient $g_t = \nabla f_t(w_t)$
\STATE Update $w_{t+1} = \Pi_{\mathcal{W}} \left( w_t - \eta_t g_t \right)$
\ENDFOR
\end{algorithmic}
\end{algorithm}

\section*{Dry Run Example}
Consider the online function $f_t(w) = (w-3)^2$ for all $t$, with $\mathcal{W} = \mathbb{R}$ (no projection needed), initial parameter $w_1 = 0$, and fixed learning rate $\eta = 0.1$. We run for $T=3$ iterations.

\textbf{Iteration 1:}
\begin{itemize}
\item Current parameter: $w_1 = 0$
\item Gradient computation: $g_1 = \nabla f_1(w_1) = 2(w_1 - 3) = 2(0 - 3) = -6$
\item Update: $w_2 = w_1 - \eta g_1 = 0 - 0.1(-6) = 0.6$
\item Objective value: $f_1(w_1) = (0-3)^2 = 9$
\end{itemize}

\textbf{Iteration 2:}
\begin{itemize}
\item Current parameter: $w_2 = 0.6$
\item Gradient computation: $g_2 = \nabla f_2(w_2) = 2(w_2 - 3) = 2(0.6 - 3) = -4.8$
\item Update: $w_3 = w_2 - \eta g_2 = 0.6 - 0.1(-4.8) = 0.6 + 0.48 = 1.08$
\item Objective value: $f_2(w_2) = (0.6-3)^2 = 5.76$
\end{itemize}

\textbf{Iteration 3:}
\begin{itemize}
\item Current parameter: $w_3 = 1.08$
\item Gradient computation: $g_3 = \nabla f_3(w_3) = 2(w_3 - 3) = 2(1.08 - 3) = -3.84$
\item Update: $w_4 = w_3 - \eta g_3 = 1.08 - 0.1(-3.84) = 1.08 + 0.384 = 1.464$
\item Objective value: $f_3(w_3) = (1.08-3)^2 = 3.6864$
\end{itemize}

\section*{Assumptions}
\begin{definition}
A function $f : \mathcal{W} \to \mathbb{R}$ is $\lambda$-strongly convex with respect to the $L_2$ norm if for all $w, v \in \mathcal{W}$ and $\alpha \in [0,1]$:
\begin{align*}
f(\alpha w + (1-\alpha)v) \leq \alpha f(w) + (1-\alpha) f(v) - \frac{\lambda}{2} \alpha(1-\alpha) \|w-v\|^2
\end{align*}
Equivalently, $f(v) \geq f(w) + \langle \nabla f(w), v - w \rangle + \frac{\lambda}{2} \|v - w\|^2$ for differentiable $f$.
\end{definition}

We assume the following for the online learning setting:
\begin{enumerate}
\item \textbf{Bounded Domain:} The feasible set $\mathcal{W} \subseteq \mathbb{R}^d$ is closed and convex, with diameter $D = \max_{w, v \in \mathcal{W}} \|w - v\| < \infty$.
\item \textbf{Strongly Convex Losses:} Each function $f_t$ is $\lambda$-strongly convex for some $\lambda > 0$.
\item \textbf{Bounded Gradients:} There exists a constant $G > 0$ such that $\|\nabla f_t(w)\| \leq G$ for all $w \in \mathcal{W}$ and all $t \in \{1, \dots, T\}$.
\end{enumerate}

\section*{Main Convergence Theorem}
\begin{theorem}
Let $R_T = \sum_{t=1}^T f_t(w_t) - \min_{w \in \mathcal{W}} \sum_{t=1}^T f_t(w)$ be the regret of Online Gradient Descent over $T$ rounds. If the loss functions $f_t$ are $\lambda$-strongly convex and the learning rate is chosen as $\eta_t = \frac{1}{\lambda t}$, then the regret satisfies:
\begin{align*}
R_T \leq \frac{G^2}{2\lambda} \sum_{t=1}^T \frac{1}{t} \leq \frac{G^2}{2\lambda} (1 + \ln T)
\end{align*}
Thus, the average regret $\frac{R_T}{T} \to 0$ as $T \to \infty$, yielding a sublinear regret of $O(\ln T)$.
\end{theorem}

\section*{Theoretical Properties}
\begin{itemize}
\item \textbf{Stability:} The diminishing step size $\eta_t = 1/(\lambda t)$ ensures that the parameter updates shrink over time, preventing oscillations and ensuring stability. The strong convexity ensures that the algorithm aggressively steps towards the optimum early on, while reducing step sizes appropriately to allow convergence of the time-varying average.
\item \textbf{Hyperparameter Sensitivity:} The optimal step size heavily depends on the strong convexity parameter $\lambda$. Underestimating $\lambda$ leads to overly conservative updates and slower convergence, while overestimating $\lambda$ leads to aggressive updates and higher instantaneous regret. Knowledge of $G$ is not strictly required for the step size, but bounds the worst-case regret.
\item \textbf{Comparison to Baselines:} For standard convex online functions, OGD with a fixed step size $\eta = O(1/\sqrt{T})$ achieves $O(\sqrt{T})$ regret. By exploiting strong convexity, we adapt the learning rate to $\eta_t = 1/(\lambda t)$, significantly improving the regret bound to $O(\ln T)$, which is optimal for this class of online optimization problems.
\end{itemize}

\section*{Discussion}
The strongly convex setting provides a more favorable geometry for online optimization, allowing the regret to grow only logarithmically with the time horizon $T$, as opposed to the square-root growth for general convex functions. This logarithmic regret implies that the time-averaged regret $\frac{R_T}{T}$ decays as $O(\frac{\ln T}{T})$, which approaches zero much faster than the $O(1/\sqrt{T})$ rate of the standard convex case. The algorithm is fully adaptive in the sense that the learning rate depends on the iteration $t$, and the proof relies on a carefully constructed potential function that telescopes over the iterations.

\section*{Preliminaries (Definitions and Lemmas)}
\begin{lemma}[Projection Lemma]
For any $u \in \mathcal{W}$ and any point $y \in \mathbb{R}^d$, let $x = \Pi_{\mathcal{W}}(y)$. Then:
\begin{align*}
\|x - u\|^2 \leq \|y - u\|^2 - \|y - x\|^2
\end{align*}
\end{lemma}
\begin{proof}
By the optimality conditions of the projection onto a convex set, for all $u \in \mathcal{W}$:
\begin{align*}
\langle y - x, u - x \rangle \leq 0
\end{align*}
Expanding $\|y - u\|^2$:
\begin{align*}
\|y - u\|^2 &= \|y - x + x - u\|^2 \\
&= \|y - x\|^2 + \|x - u\|^2 + 2 \langle y - x, x - u \rangle \\
&= \|y - x\|^2 + \|x - u\|^2 - 2 \langle y - x, u - x \rangle \\
&\geq \|y - x\|^2 + \|x - u\|^2
\end{align*}
Rearranging yields the stated bound.
\end{proof}

\begin{lemma}[Strong Convexity Inequality]
If $f_t$ is $\lambda$-strongly convex, then for all $w, v \in \mathcal{W}$:
\begin{align*}
f_t(w) - f_t(v) \leq \langle \nabla f_t(w), w - v \rangle - \frac{\lambda}{2} \|w - v\|^2
\end{align*}
\end{lemma}
\begin{proof}
This follows directly from the definition of $\lambda$-strong convexity applied to the differentiable function $f_t$, which guarantees:
\begin{align*}
f_t(v) \geq f_t(w) + \langle \nabla f_t(w), v - w \rangle + \frac{\lambda}{2} \|v - w\|^2
\end{align*}
Subtracting $f_t(v)$ from both sides and multiplying by $-1$ yields the result.
\end{proof}

\section*{Main Proof (Step-by-Step Derivation)}
Let $w^* = \arg\min_{w \in \mathcal{W}} \sum_{t=1}^T f_t(w)$ be the optimal offline comparator. We define the potential function $\Phi_t = \frac{1}{2\eta_t} \|w_t - w^*\|^2$. 

[Step 1] We evaluate the difference in potentials between consecutive steps:
\begin{align*}
\Phi_t - \Phi_{t+1} = \frac{1}{2\eta_t} \|w_t - w^*\|^2 - \frac{1}{2\eta_{t+1}} \|w_{t+1} - w^*\|^2
\end{align*}

[Step 2] Substitute the learning rate schedule $\eta_t = \frac{1}{\lambda t}$ and $\eta_{t+1} = \frac{1}{\lambda(t+1)}$:
\begin{align*}
\Phi_t - \Phi_{t+1} = \frac{\lambda t}{2} \|w_t - w^*\|^2 - \frac{\lambda (t+1)}{2} \|w_{t+1} - w^*\|^2
\end{align*}

[Step 3] Apply the Projection Lemma (Lemma 1) with $y = w_t - \eta_t \nabla f_t(w_t)$, $x = w_{t+1}$, and $u = w^*$:
\begin{align*}
\|w_{t+1} - w^*\|^2 \leq \|w_t - \eta_t \nabla f_t(w_t) - w^*\|^2 - \|w_t - \eta_t \nabla f_t(w_t) - w_{t+1}\|^2
\end{align*}
Since the second term is non-negative, we drop it to establish an inequality:
\begin{align*}
\|w_{t+1} - w^*\|^2 \leq \|w_t - w^*\|^2 - 2\eta_t \langle \nabla f_t(w_t), w_t - w^* \rangle + \eta_t^2 \|\nabla f_t(w_t)\|^2
\end{align*}

[Step 4] Multiply the inequality in [Step 3] by $-\frac{\lambda(t+1)}{2}$ and add $\frac{\lambda t}{2} \|w_t - w^*\|^2$:
\begin{align*}
\Phi_t - \Phi_{t+1} &\geq \frac{\lambda t}{2} \|w_t - w^*\|^2 - \frac{\lambda (t+1)}{2} \left( \|w_t - w^*\|^2 - 2\eta_t \langle \nabla f_t(w_t), w_t - w^* \rangle + \eta_t^2 \|\nabla f_t(w_t)\|^2 \right) \\
&= \frac{\lambda t}{2} \|w_t - w^*\|^2 - \frac{\lambda t}{2} \|w_t - w^*\|^2 - \frac{\lambda}{2} \|w_t - w^*\|^2 \\
&\quad + \lambda(t+1)\eta_t \langle \nabla f_t(w_t), w_t - w^* \rangle - \frac{\lambda(t+1)\eta_t^2}{2} \|\nabla f_t(w_t)\|^2
\end{align*}

[Step 5] Simplify the expression by canceling the $\frac{\lambda t}{2} \|w_t - w^*\|^2$ terms:
\begin{align*}
\Phi_t - \Phi_{t+1} \geq - \frac{\lambda}{2} \|w_t - w^*\|^2 + \lambda(t+1)\eta_t \langle \nabla f_t(w_t), w_t - w^* \rangle - \frac{\lambda(t+1)\eta_t^2}{2} \|\nabla f_t(w_t)\|^2
\end{align*}

[Step 6] Substitute $\eta_t = \frac{1}{\lambda t}$ into the remaining terms:
\begin{align*}
\lambda(t+1)\eta_t &= \lambda(t+1) \frac{1}{\lambda t} = \frac{t+1}{t} = 1 + \frac{1}{t} \\
\frac{\lambda(t+1)\eta_t^2}{2} &= \frac{\lambda(t+1)}{2} \frac{1}{\lambda^2 t^2} = \frac{t+1}{2\lambda t^2}
\end{align*}
Thus, the inequality becomes:
\begin{align*}
\Phi_t - \Phi_{t+1} \geq - \frac{\lambda}{2} \|w_t - w^*\|^2 + \left(1 + \frac{1}{t}\right) \langle \nabla f_t(w_t), w_t - w^* \rangle - \frac{t+1}{2\lambda t^2} \|\nabla f_t(w_t)\|^2
\end{align*}

[Step 7] Rearrange the above inequality to isolate $\langle \nabla f_t(w_t), w_t - w^* \rangle$:
\begin{align*}
\left(1 + \frac{1}{t}\right) \langle \nabla f_t(w_t), w_t - w^* \rangle \leq \Phi_t - \Phi_{t+1} + \frac{\lambda}{2} \|w_t - w^*\|^2 + \frac{t+1}{2\lambda t^2} \|\nabla f_t(w_t)\|^2
\end{align*}
Multiplying both sides by $\frac{t}{t+1}$:
\begin{align*}
\langle \nabla f_t(w_t), w_t - w^* \rangle \leq \frac{t}{t+1} \left( \Phi_t - \Phi_{t+1} \right) + \frac{t}{t+1} \frac{\lambda}{2} \|w_t - w^*\|^2 + \frac{1}{2\lambda t} \|\nabla f_t(w_t)\|^2
\end{align*}

[Step 8] Apply Lemma 2 (Strong Convexity Inequality) to bound the instantaneous regret $f_t(w_t) - f_t(w^*)$:
\begin{align*}
f_t(w_t) - f_t(w^*) \leq \langle \nabla f_t(w_t), w_t - w^* \rangle - \frac{\lambda}{2} \|w_t - w^*\|^2
\end{align*}
Substituting the result from [Step 7]:
\begin{align*}
f_t(w_t) - f_t(w^*) &\leq \frac{t}{t+1} \left( \Phi_t - \Phi_{t+1} \right) + \frac{t}{t+1} \frac{\lambda}{2} \|w_t - w^*\|^2 - \frac{\lambda}{2} \|w_t - w^*\|^2 + \frac{1}{2\lambda t} \|\nabla f_t(w_t)\|^2
\end{align*}

[Step 9] Combine the terms involving $\frac{\lambda}{2} \|w_t - w^*\|^2$:
\begin{align*}
\frac{t}{t+1} \frac{\lambda}{2} \|w_t - w^*\|^2 - \frac{\lambda}{2} \|w_t - w^*\|^2 = \left( \frac{t}{t+1} - 1 \right) \frac{\lambda}{2} \|w_t - w^*\|^2 = -\frac{1}{t+1} \frac{\lambda}{2} \|w_t - w^*\|^2
\end{align*}
Since $-\frac{1}{t+1} \frac{\lambda}{2} \|w_t - w^*\|^2 \leq 0$, we can safely drop this term to establish an upper bound:
\begin{align*}
f_t(w_t) - f_t(w^*) \leq \frac{t}{t+1} \left( \Phi_t - \Phi_{t+1} \right) + \frac{1}{2\lambda t} \|\nabla f_t(w_t)\|^2
\end{align*}

[Step 10] Sum the instantaneous regret from $t=1$ to $T$:
\begin{align*}
\sum_{t=1}^T \left( f_t(w_t) - f_t(w^*) \right) \leq \sum_{t=1}^T \frac{t}{t+1} \left( \Phi_t - \Phi_{t+1} \right) + \sum_{t=1}^T \frac{1}{2\lambda t} \|\nabla f_t(w_t)\|^2
\end{align*}

[Step 11] Expand the first summation on the right-hand side:
\begin{align*}
\sum_{t=1}^T \frac{t}{t+1} \left( \Phi_t - \Phi_{t+1} \right) &= \frac{1}{2}(\Phi_1 - \Phi_2) + \frac{2}{3}(\Phi_2 - \Phi_3) + \dots + \frac{T}{T+1}(\Phi_T - \Phi_{T+1}) \\
&= \frac{1}{2}\Phi_1 + \left(\frac{2}{3} - \frac{1}{2}\right)\Phi_2 + \left(\frac{3}{4} - \frac{2}{3}\right)\Phi_3 + \dots + \left(\frac{T}{T+1} - \frac{T-1}{T}\right)\Phi_T - \frac{T}{T+1}\Phi_{T+1}
\end{align*}

[Step 12] Compute the coefficient for the generic term $\Phi_t$ where $2 \leq t \leq T$:
\begin{align*}
\frac{t}{t+1} - \frac{t-1}{t} = \frac{t^2 - (t^2 - 1)}{t(t+1)} = \frac{1}{t(t+1)}
\end{align*}
Thus, the sum becomes:
\begin{align*}
\sum_{t=1}^T \frac{t}{t+1} \left( \Phi_t - \Phi_{t+1} \right) = \frac{1}{2}\Phi_1 + \sum_{t=2}^T \frac{1}{t(t+1)} \Phi_t - \frac{T}{T+1}\Phi_{T+1}
\end{align*}
Since $\Phi_t \geq 0$ and $\frac{T}{T+1}\Phi_{T+1} \geq 0$, we can drop the last term to get an upper bound:
\begin{align*}
\sum_{t=1}^T \frac{t}{t+1} \left( \Phi_t - \Phi_{t+1} \right) \leq \frac{1}{2}\Phi_1 + \sum_{t=2}^T \frac{1}{t(t+1)} \Phi_t
\end{align*}

[Step 13] Substitute the definition $\Phi_t = \frac{\lambda t}{2} \|w_t - w^*\|^2$ back into the bound:
\begin{align*}
\frac{1}{2}\Phi_1 &= \frac{\lambda}{2} \|w_1 - w^*\|^2 \leq \frac{\lambda D^2}{2} \\
\frac{1}{t(t+1)} \Phi_t &= \frac{1}{t(t+1)} \frac{\lambda t}{2} \|w_t - w^*\|^2 = \frac{\lambda}{2(t+1)} \|w_t - w^*\|^2 \leq \frac{\lambda D^2}{2(t+1)}
\end{align*}

[Step 14] Bound the telescoping series using the domain diameter $D$:
\begin{align*}
\frac{1}{2}\Phi_1 + \sum_{t=2}^T \frac{1}{t(t+1)} \Phi_t \leq \frac{\lambda D^2}{2} + \sum_{t=2}^T \frac{\lambda D^2}{2(t+1)} = \frac{\lambda D^2}{2} \left( 1 + \sum_{t=2}^T \frac{1}{t+1} \right) = \frac{\lambda D^2}{2} \sum_{t=1}^T \frac{1}{t+1}
\end{align*}

[Step 15] Apply the bound on gradients $\|\nabla f_t(w_t)\|^2 \leq G^2$ to the second summation from [Step 10]:
\begin{align*}
\sum_{t=1}^T \frac{1}{2\lambda t} \|\nabla f_t(w_t)\|^2 \leq \sum_{t=1}^T \frac{G^2}{2\lambda t}
\end{align*}

[Step 16] Combine the results from [Step 10], [Step 14], and [Step 15] to bound the total regret $R_T$:
\begin{align*}
R_T = \sum_{t=1}^T \left( f_t(w_t) - f_t(w^*) \right) \leq \frac{\lambda D^2}{2} \sum_{t=1}^T \frac{1}{t+1} + \frac{G^2}{2\lambda} \sum_{t=1}^T \frac{1}{t}
\end{align*}

\section*{Rate Derivation (With Constants)}
From [Step 16], we have:
\begin{align*}
R_T \leq \frac{G^2}{2\lambda} \sum_{t=1}^T \frac{1}{t} + \frac{\lambda D^2}{2} \sum_{t=1}^T \frac{1}{t+1}
\end{align*}
We use the standard bound for the Harmonic series: $\sum_{t=1}^T \frac{1}{t} \leq 1 + \ln T$ and $\sum_{t=1}^T \frac{1}{t+1} \leq \ln(T+1) \leq \ln T + 1$. Applying these bounds:
\begin{align*}
R_T &\leq \frac{G^2}{2\lambda} (1 + \ln T) + \frac{\lambda D^2}{2} (1 + \ln T) \\
&= \left( \frac{G^2 + \lambda^2 D^2}{2\lambda} \right) (1 + \ln T)
\end{align*}
This confirms that the total regret scales as $R_T = O(\ln T)$, and the average regret scales as $\frac{R_T}{T} = O\left(\frac{\ln T}{T}\right)$, which vanishes as $T \to \infty$. The constant multiplier is $\frac{G^2 + \lambda^2 D^2}{2\lambda}$.

\section*{Special Cases}
\textbf{Case 1: Fixed Learning Rate for General Convex Functions.}
If the online functions are convex but not strongly convex ($\lambda = 0$), the learning rate $\eta_t = \frac{1}{\lambda t}$ is undefined. In this case, the standard OGD analysis dictates a fixed or decreasing step size $\eta_t = \frac{D}{G\sqrt{t}}$. By modifying the potential to $\Phi_t = \frac{1}{2\eta} \|w_t - w^*\|^2$ and summing over $T$ steps, the regret evaluates to:
\begin{align*}
R_T \leq \frac{DG}{2} + \frac{DG}{2} \sum_{t=2}^T \frac{1}{\sqrt{t}} \leq \frac{3}{2} DG \sqrt{T}
\end{align*}
yielding the well-known $O(\sqrt{T})$ sublinear regret bound for general convex online functions.

\textbf{Case 2: Known Time Horizon $T$.}
If the total number of iterations $T$ is known a priori, we can use a constant step size $\eta = \frac{D}{G\sqrt{T}}$ for general convex functions. This leads to a simplified proof where the sum of squared norms is bounded by $T G^2$, resulting exactly in $R_T \leq DG\sqrt{T}$. For the strongly convex case, knowing $T$ does not improve the asymptotic rate of $O(\ln T)$, but allows slight optimizations in the additive constants of the learning rate $\eta_t = \frac{1}{\lambda (t + c)}$ where $c$ is calibrated based on $T$.

\end{document}